\section{Maximum XOR of Two Numbers in an Array}
\label{sec:trie-max-xor}

\problemstatement{Given an array $\textit{nums}$ of non-negative
integers, find the maximum value of $\textit{nums}[i] \oplus
\textit{nums}[j]$ over all pairs $i \ne j$ (where $\oplus$ denotes
bitwise XOR).}

\begin{patternbox}
\emph{``Maximum XOR''} between pairs of numbers is the keyword for a
\textbf{binary trie}: treat each number as a fixed-length sequence of
bits (most significant bit first) and insert it into a trie exactly as
Section~\ref{sec:trie-implement} inserts characters of a string -- except
the ``alphabet'' here has only two symbols, $0$ and $1$. Once every
number is in the trie, maximizing the XOR with a given number becomes a
greedy walk: at each bit position, XOR is maximized by making that bit
\textbf{differ}, so we always prefer the branch \emph{opposite} to the
current number's bit, whenever it exists.
\end{patternbox}

\subsection*{Understanding the problem carefully}

XOR-ing two bits produces $1$ exactly when they differ. To make the
overall XOR as large as possible, we want to maximize the \emph{most
significant} differing bit first (since higher bit positions dominate
the numeric value), then the next, and so on -- exactly the same
most-significant-first greedy logic used throughout binary search's
``binary search on the answer'' problems, but here applied bit by bit
through a trie rather than narrowing a numeric range.

\begin{ideabox}[Insert All Numbers as Bit-Paths, Then Greedily Walk for the Opposite Bit]
Insert every number into a binary trie, one bit at a time from the most
significant bit downward (using a fixed bit width, e.g.\ $32$, so all
numbers have aligned representations). For each number $x$ in the array,
walk the trie trying, at each bit position, to move into the child
representing the \emph{opposite} of $x$'s bit at that position; if that
child does not exist, settle for the only available child (the same
bit). Accumulate the resulting XOR value bit by bit as you walk, and
track the maximum such value found over all numbers $x$.
\end{ideabox}

\subsection*{Why the greedy opposite-bit choice is always optimal}

Because bit positions contribute to the final value with strictly
decreasing weight (the most significant bit contributes more than the
sum of all lower bits combined, for a fixed bit width), it is always
better to secure a $1$ at a higher bit position than to optimize any
combination of lower bits -- so greedily taking the opposite-bit branch
whenever it is available, at every level from the most significant bit
down, is guaranteed to produce the maximum possible XOR with the
existing numbers in the trie, by the same most-significant-digit-first
reasoning used whenever comparing fixed-width binary numbers.

\subsection*{Algorithm}

\begin{enumerate}
  \item Insert every number in $\textit{nums}$ into a binary trie, bit
        by bit from the most significant bit (say bit $31$) to the
        least significant (bit $0$).
  \item For each number $x$ in $\textit{nums}$:
  \begin{enumerate}
    \item Walk the trie from the root. At each bit position $b$ (from
          most to least significant): let $\textit{desired}$ be the
          opposite of $x$'s bit $b$. If the trie node has a child for
          $\textit{desired}$, move into it and set this bit of the
          running XOR result to $1$; otherwise, move into the only
          available child and set this bit to $0$.
    \item Update the global maximum with this number's resulting XOR
          value.
  \end{enumerate}
  \item Return the global maximum.
\end{enumerate}

\subsection*{Dry run}

\dryrun{Take $\textit{nums} = [3,10,5,25,2,8]$, using $5$-bit
representations for clarity ($3{=}00011$, $10{=}01010$, $5{=}00101$,
$25{=}11001$, $2{=}00010$, $8{=}01000$). Query the best XOR partner for
$5{=}00101$.}

\begin{itemize}
  \item Bit $4$ (value $16$) of $5$ is $0$; desired opposite is $1$. Only
        $25$ has bit $4=1$. Move there; XOR bit $4=1$.
  \item Bit $3$ (value $8$) of $5$ is $0$; desired $1$. Within this
        branch, only $25$ remains, and its bit $3=1$ -- matches desired.
        Move there; XOR bit $3=1$.
  \item Bit $2$ (value $4$) of $5$ is $1$; desired $0$. $25$'s bit
        $2=0$ -- matches. Move there; XOR bit $2=1$.
  \item Bit $1$ (value $2$) of $5$ is $0$; desired $1$. Only $25$'s
        branch remains here, and its bit $1=0$ -- not the desired
        branch, but it's the only one available. Move there; XOR bit
        $1=0$.
  \item Bit $0$ (value $1$) of $5$ is $1$; desired $0$. Only $25$'s
        bit $0=1$ available. Move there; XOR bit $0=0$.
\end{itemize}

Resulting XOR bits: $11100_2 = 28$. Indeed $5 \oplus 25 = 28$, the
maximum pairwise XOR for this array.

\subsection*{C++ implementation}

\begin{lstlisting}
#include <bits/stdc++.h>
using namespace std;

struct TrieNode {
    TrieNode* children[2] = {nullptr, nullptr};
};

class BinaryTrie {
    TrieNode* root = new TrieNode();
    static const int BITS = 32;

public:
    void insert(int num) {
        TrieNode* node = root;
        for (int b = BITS - 1; b >= 0; b--) {
            int bit = (num >> b) & 1;
            if (!node->children[bit]) node->children[bit] = new TrieNode();
            node = node->children[bit];
        }
    }

    int maxXorWith(int num) {
        TrieNode* node = root;
        int result = 0;
        for (int b = BITS - 1; b >= 0; b--) {
            int bit = (num >> b) & 1;
            int desired = 1 - bit;
            if (node->children[desired]) {
                result |= (1 << b);
                node = node->children[desired];
            } else {
                node = node->children[bit];
            }
        }
        return result;
    }
};

int findMaximumXOR(vector<int>& nums) {
    BinaryTrie trie;
    for (int x : nums) trie.insert(x);

    int maxXor = 0;
    for (int x : nums) {
        maxXor = max(maxXor, trie.maxXorWith(x));
    }
    return maxXor;
}

int main() {
    vector<int> nums = {3, 10, 5, 25, 2, 8};
    cout << "Maximum XOR: " << findMaximumXOR(nums) << endl;
    return 0;
}
\end{lstlisting}

\begin{complexitybox}
\textbf{Time Complexity.} Inserting all $n$ numbers costs $O(n \cdot
\textit{BITS})$. Querying the best partner for each of the $n$ numbers
costs another $O(n \cdot \textit{BITS})$. Overall:
\[
  O(n \cdot \textit{BITS}) = O(n),
\]
treating the fixed bit width ($32$) as a constant -- a dramatic
improvement over the $O(n^2)$ brute-force check of every pair.
\textbf{Space Complexity.} $O(n \cdot \textit{BITS})$ for the trie in the
worst case.
\end{complexitybox}

\begin{takeawaybox}
A binary trie is a character trie with a two-letter alphabet -- every
technique from Sections~\ref{sec:trie-implement}--\ref{sec:trie-distinct-substrings}
carries over unchanged, just with bits instead of letters. Whenever a
problem involves maximizing or counting based on XOR (or, more generally,
bitwise structure) across many numbers, a binary trie processed
most-significant-bit-first is almost always the intended technique.
\end{takeawaybox}
